% Bagian: computability
% Bab: computability-theory
% Subbagian: rice-theorem

\documentclass[../../../include/open-logic-section]{subfiles}

\begin{document}

\olfileid[id]{cmp}{thy}{rce}
\olsection{Teorema Rice}

Jika dipikirkan, Anda akan melihat bahwa kekhususan $\fn{Tot}$ tidak berperan
dalam bukti \olref[tot]{prop:total}. Kita merancang $h(x,y)$ agar berperilaku
seperti fungsi konstan $j(y)=0$ tepat ketika $x$ berada dalam~$K$; tetapi dalam
keadaan itu kita dapat juga membuatnya berperilaku seperti sembarang fungsi
dapat dihitung parsial lainnya. Pengamatan ini memungkinkan kita menyatakan
teorema yang lebih umum, yang secara kasar mengatakan bahwa tidak ada sifat
nontrivial dari fungsi dapat dihitung yang dapat diputuskan.

Ingatlah bahwa $\cfind{0}$, $\cfind{1}$, $\cfind{2}$,~\dots merupakan
enumerasi standar kita atas fungsi-fungsi dapat dihitung parsial.

\begin{thm}[Teorema Rice]
  Misalkan $C$ sembarang himpunan fungsi dapat dihitung parsial dan misalkan
  $A = \Setabs{n}{\cfind{n} \in C}$. Jika $A$ dapat dihitung, maka $C$ kosong
  atau $C$ merupakan himpunan semua fungsi dapat dihitung parsial.
\end{thm}

Suatu \emph{himpunan indeks} adalah himpunan $A$ dengan sifat bahwa jika $n$
dan $m$ merupakan indeks yang ``menghitung'' fungsi yang sama, maka $n$ dan
$m$ keduanya berada dalam~$A$, atau keduanya tidak berada di dalamnya. Tidak
sulit melihat bahwa himpunan~$A$ dalam teorema mempunyai sifat ini. Sebaliknya,
jika $A$ merupakan himpunan indeks dan $C$ merupakan himpunan fungsi yang
dihitung oleh indeks-indeks tersebut, maka
$A = \Setabs{n}{\cfind{n} \in C}$.

\begin{explain}
Dengan terminologi ini, teorema Rice ekuivalen dengan pernyataan bahwa tidak
ada himpunan indeks nontrivial yang dapat diputuskan. Untuk memahami isi
teorema tersebut, akan membantu jika kita menekankan perbedaan antara
\emph{program} (misalnya dalam bahasa pemrograman kesukaan Anda) dan fungsi
yang dihitungnya. Tentu ada pertanyaan tentang program (indeks), yang
merupakan objek sintaktis, yang dapat diputuskan secara komputasional: apakah
program ini memuat lebih dari 150 simbol? Apakah program ini mempunyai lebih
dari 22 baris? Apakah program ini mempunyai pernyataan ``while''? Apakah
untai ``hello world'' pernah muncul sebagai argumen bagi pernyataan
``print''? Teorema Rice mengatakan bahwa tidak ada pertanyaan nontrivial
tentang \emph{perilaku} program yang dapat diputuskan secara komputasional.
Ini mencakup pertanyaan seperti: apakah program berhenti pada masukan~$0$?
Apakah program pernah berhenti? Apakah program pernah mengeluarkan bilangan
genap?
\end{explain}

\begin{proof}[Bukti teorema Rice]
Andaikan $C$ tidak kosong dan bukan pula himpunan semua fungsi dapat dihitung
parsial, dan misalkan $A$ himpunan indeks fungsi-fungsi dalam~$C$. Kita akan
menunjukkan bahwa jika $A$ dapat dihitung, kita dapat menyelesaikan masalah
penghentian; jadi $A$ tidak dapat dihitung.

Tanpa mengurangi keumuman, kita dapat mengandaikan bahwa fungsi $f$ yang tidak
terdefinisi di mana pun tidak berada dalam~$C$ (jika tidak, tukarkan $C$
dengan komplemennya dalam argumen di bawah). Misalkan $g$ sembarang fungsi
dalam~$C$. Gagasannya ialah bahwa jika kita dapat memutuskan~$A$, kita dapat
membedakan indeks yang menghitung~$f$ dari indeks yang menghitung~$g$; lalu
kita dapat menggunakan kemampuan itu untuk menyelesaikan masalah penghentian.

Caranya sebagai berikut. Dengan menggunakan fungsi dapat dihitung parsial
universal, kita dapat mendefinisikan fungsi
\[
h(x,y) \simeq
\begin{cases}
\text{tidak terdefinisi} & \text{jika $\cfind{x}(x) \fundefined$} \\
g(y) & \text{selain itu.}
\end{cases}
\]
Untuk menghitung $h$, mula-mula kita mencoba menghitung $\cfind{x}(x)$; jika
komputasi itu berhenti, kita melanjutkan dengan menghitung~$g(y)$; dan jika
\emph{komputasi itu} berhenti, kita mengembalikan keluarannya. Secara lebih
formal, kita dapat menulis
\[
h(x,y) \simeq \Proj{2}{0}(g(y),\fn{Un}(x,x)).
\]
Di sini $\Proj{2}{0}(z_0,z_1)=z_0$ adalah fungsi proyeksi berargumen~2 yang
mengembalikan argumen ke-$0$, dan fungsi ini dapat dihitung.

Dengan demikian, $h$ merupakan komposisi fungsi-fungsi dapat dihitung parsial,
dan ruas kanan terdefinisi serta sama dengan~$g(y)$ tepat ketika
$\fn{Un}(x,x)$ dan $g(y)$ keduanya terdefinisi.

Perhatikan bahwa untuk $x$ tetap, jika $\cfind{x}(x)$ tidak terdefinisi, maka
$h(x,y)$ tidak terdefinisi bagi setiap~$y$; dan jika $\cfind{x}(x)$
terdefinisi, maka $h(x,y) \simeq g(y)$. Jadi, bagi setiap nilai tetap~$x$,
$h(x,y)$ berperilaku tepat seperti $f$ atau tepat seperti $g$, dan memutuskan
apakah $\cfind{x}(x)$ terdefinisi sama dengan memutuskan kasus mana yang
berlaku. Namun, ini sama dengan memutuskan apakah $h_x(y) \simeq h(x,y)$
berada dalam~$C$ atau tidak; jika $A$ dapat dihitung, kita dapat melakukan
hal tersebut.

Secara lebih formal, karena $h$ dapat dihitung parsial, $h$ sama dengan fungsi
$\cfind{e}$ bagi suatu indeks~$e$. Menurut teorema $s$-$m$-$n$, terdapat
fungsi rekursif primitif~$s$ sedemikian sehingga bagi setiap~$x$,
$\cfind{s(e,x)}(y)\simeq h_x(y)$. Sekarang, bagi setiap $x$, jika
$\cfind{x}(x) \fdefined$, maka $\cfind{s(e,x)}$ merupakan fungsi yang sama
dengan~$g$, sehingga $s(e,x)$ berada dalam~$A$. Di lain pihak, jika
$\cfind{x}(x) \uparrow$, maka $\cfind{s(e,x)}$ merupakan fungsi yang sama
dengan~$f$, sehingga $s(e,x)$ tidak berada dalam~$A$. Dengan kata lain, bagi
setiap~$x$, $x \in K$ jika dan hanya jika $s(e,x) \in A$. Jika $A$ dapat
dihitung, $K$ juga dapat dihitung, yang merupakan suatu kontradiksi. Jadi,
$A$ tidak dapat dihitung.
\end{proof}

Teorema Rice sangat kuat. Akibat langsung berikut menunjukkan beberapa
contoh penerapannya.
\begin{cor}
Himpunan-himpunan berikut tidak dapat diputuskan.
\begin{enumerate}
\item $\Setabs{x}{\text{$17$ berada dalam daerah hasil $\cfind{x}$}}$
\item $\Setabs{x}{\text{$\cfind{x}$ konstan}}$
\item $\Setabs{x}{\text{$\cfind{x}$ total}}$
\item $\Setabs{x}{\text{setiap kali $y < y'$, $\cfind{x}(y) \fdefined$, dan
    jika $\cfind{x}(y') \fdefined$, maka $\cfind{x}(y) < \cfind{x}(y')$}}$
\end{enumerate}
\end{cor}

\begin{proof}
  Semuanya merupakan himpunan indeks nontrivial.
\end{proof}

\end{document}
